%%%
%%% Radical "⾜"
%%%
\section*{Radical 157: ``⾜''}\addcontentsline{toc}{section}{Radical 157: ⾜}\addcontentsline{loh}{figure}{\#\#\#\# 157: ⾜}

%%%%%%%%%% 足 %%%%%%%%%%
\subsection*{足}\addcontentsline{loh}{figure}{足}

\begin{Entry}{足}{7}{⾜}[Kangxi 157]
  \begin{Phonetics}{足}{ju4}
    \definition{adj.}{excessivo}
  \end{Phonetics}
  \begin{Phonetics}{足}{zu2}[][HSK 6]
    \definition{adj.}{amplo; bastante; suficiente}
    \definition{adv.}{totalmente; tanto quanto; indica suficiente até certo ponto ou grau | (geralmente no negativo) o suficiente; suficientemente; é totalmente possível; vale a pena}
    \definition{s.}{pé; um termo geral para os membros inferiores do corpo humano; especificamente a parte abaixo do tornozelo | futebol; refere"-se ao futebol ou a um time de futebol | pé (de utensílios); a parte inferior do objeto tem o formato de uma pé e serve de suporte}
  \seealsoref{够}{gou4}
  \synonymref{脚}{jiao3}
  \synonymref{腿}{tui3}
  \antonymref{亏}{kui1}
  \end{Phonetics}
\end{Entry}

\begin{Entry}{足以}{7,4}{⾜,⼈}
  \begin{Phonetics}{足以}{zu2yi3}[][HSK 6]
    \definition{adj.}{suficiente; totalmente capaz de; o suficiente para fazer algo}
    \definition{v.}{bastar}
  \synonymref{得以}{de2yi3}
  \synonymref{可以}{ke3yi3}
  \synonymref{能够}{neng2gou4}
  \synonymref{知足}{zhi1zu2}
  \synonymref{足够}{zu2gou4}
  \antonymref{不够}{bu2gou4}
  \antonymref{不足}{bu4zu2}
  \antonymref{难以}{nan2yi3}
  \end{Phonetics}
\end{Entry}

\begin{Entry}{足月}{7,4}{⾜,⽉}
  \begin{Phonetics}{足月}{zu2 yue4}
    \definition{adj.}{(um feto) completo; maduro | completa (gestação)}
  \end{Phonetics}
\end{Entry}

\begin{Entry}{足足}{7,7}{⾜,⾜}
  \begin{Phonetics}{足足}{zu2zu2}
    \definition{adv.}{tanto quanto; totalmente; indica que o número não é menor ou é completamente suficiente}
  \synonymref{整整}{zheng3zheng3}
  \antonymref{缺乏}{que1fa2}
  \end{Phonetics}
\end{Entry}

\begin{Entry}{足迹}{7,9}{⾜,⾡}
  \begin{Phonetics}{足迹}{zu2ji4}[][HSK 7-9]
    \definition[个]{s.}{rastro; pegada; marcas de pegadas; referindo"-se metaforicamente a lugares visitados}
  \synonymref{脚印}{jiao3yin4}
  \end{Phonetics}
\end{Entry}

\begin{Entry}{足够}{7,11}{⾜,⼣}
  \begin{Phonetics}{足够}{zu2gou4}[][HSK 3]
    \definition{adj.}{bastante; amplo; suficiente; atingir o nível adequado ou capaz de satisfazer as necessidades}
    \definition{v.}{satisfazer; ser suficiente; estar a contento}
  \synonymref{充分}{chong1fen4}
  \synonymref{充满}{chong1man3}
  \synonymref{充沛}{chong1pei4}
  \synonymref{充实}{chong1shi2}
  \synonymref{充足}{chong1zu2}
  \synonymref{丰富}{feng1fu4}
  \synonymref{足以}{zu2yi3}
  \antonymref{不足}{bu4zu2}
  \antonymref{欠缺}{qian4que1}
  \antonymref{缺乏}{que1fa2}
  \end{Phonetics}
\end{Entry}

\begin{Entry}{足球}{7,11}{⾜,⽟}
  \begin{Phonetics}{足球}{zu2qiu2}[][HSK 3]
    \definition[个,只,颗,袋]{s.}{futebol | bola de futebol}
  \end{Phonetics}
\end{Entry}

\begin{Entry}{足球队}{7,11,4}{⾜,⽟,⾩}
  \begin{Phonetics}{足球队}{zu2qiu2dui4}
    \definition[个,支]{s.}{time de futebol}
  \end{Phonetics}
\end{Entry}

\begin{Entry}{足球协会}{7,11,6,6}{⾜,⽟,⼗,⼈}
  \begin{Phonetics}{足球协会}{zu2qiu2 xie2hui4}
    \definition*{s.}{Associação de Futebol | Liga de Futebol}
  \end{Phonetics}
\end{Entry}

\begin{Entry}{足球场}{7,11,6}{⾜,⽟,⼟}
  \begin{Phonetics}{足球场}{zu2qiu2chang3}
    \definition[个,座]{s.}{campo de futebol; campo de futebol americano; o campo usado para jogar futebol}
  \end{Phonetics}
\end{Entry}

\begin{Entry}{足球迷}{7,11,9}{⾜,⽟,⾡}
  \begin{Phonetics}{足球迷}{zu2qiu2mi2}
    \definition{s.}{fã (ou entusiasta) de futebol}
  \end{Phonetics}
\end{Entry}

\begin{Entry}{足球赛}{7,11,14}{⾜,⽟,⾙}
  \begin{Phonetics}{足球赛}{zu2qiu2sai4}
    \definition[场]{s.}{partida (ou jogo) de futebol | competição de futebol}
  \end{Phonetics}
\end{Entry}

\begin{Entry}{足智多谋}{7,12,6,11}{⾜,⽇,⼣,⾔}
  \begin{Phonetics}{足智多谋}{zu2zhi4-duo1mou2}[][HSK 7-9]
    \definition{expr.}{sábio e cheio de estratagemas; sábio e engenhoso; cheio de estratagemas; engenhoso}
  \antonymref{一筹莫展}{yi4chou2-mo4zhan3}
  \end{Phonetics}
\end{Entry}

%%%%%%%%%% 趴 %%%%%%%%%%
\subsection*{趴}\addcontentsline{loh}{figure}{趴}

\begin{Entry}{趴}{9}{⾜}
  \begin{Phonetics}{趴}{pa1}[][HSK 7-9]
    \definition{v.}{deitar"-se de bruços; arriar; espreguiçar"-se | curvar"-se; apoiar"-se em; inclinar"-se para a frente apoiando"-se em um objeto}
  \end{Phonetics}
\end{Entry}

%%%%%%%%%% 距 %%%%%%%%%%
\subsection*{距}\addcontentsline{loh}{figure}{距}

\begin{Entry}{距}{11}{⾜}
  \begin{Phonetics}{距}{ju4}[][HSK 7-9]
    \definition{s.}{distância | espora (de um galo, etc.)}
    \definition{v.}{estar separado (longe) de; estar distante de}
  \end{Phonetics}
\end{Entry}

\begin{Entry}{距离}{11,10}{⾜,⼇}
  \begin{Phonetics}{距离}{ju4li2}[][HSK 4]
    \definition[个,段]{s.}{distância}
    \definition{v.}{estar distante de}
  \seealsoref{距}{ju4}
  \seealsoref{离}{li2}
  \synonymref{差距}{cha1ju4}
  \synonymref{间隔}{jian4ge2}
  \end{Phonetics}
\end{Entry}

%%%%%%%%%% 跌 %%%%%%%%%%
\subsection*{跌}\addcontentsline{loh}{figure}{跌}

\begin{Entry}{跌}{12}{⾜}
  \begin{Phonetics}{跌}{die1}[][HSK 6]
    \definition{s.}{(de um objeto, etc.) queda; tombo | (de preços, etc.) queda}
    \definition{v.}{cair; tombar; perder o equilíbrio e cair | cair (objetos caindo); descer | cair (queda de preços)}
  \synonymref{降}{jiang4}
  \synonymref{坠}{zhui4}
  \antonymref{升}{sheng1}
  \antonymref{涨}{zhang3}
  \end{Phonetics}
\end{Entry}

%%%%%%%%%% 跑 %%%%%%%%%%
\subsection*{跑}\addcontentsline{loh}{figure}{跑}

\begin{Entry}{跑}{12}{⾜}
  \begin{Phonetics}{跑}{pao2}
    \definition{v.}{(animais) bater com a pata (no chão); (animais) escavar o solo com suas garras ou cascos}
  \end{Phonetics}
  \begin{Phonetics}{跑}{pao3}[][HSK 1]
    \definition{v.}{correr; pessoas ou animais que se movem rapidamente para a frente com as pernas e os pés | caminhar; passear | fugir; escapar | correr de um lado para outro; fazer rondas; correr atrás de algo | de um líquido ou gás) vazar; evaporar | (como complemento de um verbo) fora; longe | participar de uma corrida}
  \end{Phonetics}
\end{Entry}

\begin{Entry}{跑马}{12,3}{⾜,⾺}
  \begin{Phonetics}{跑马}{pao3ma3}
    \definition{s.}{corrida de cavalos | cavalo de corrida}
    \definition{v.}{andar a cavalo; dar uma volta a cavalo; correr a cavalo}
  \end{Phonetics}
\end{Entry}

\begin{Entry}{跑车}{12,4}{⾜,⾞}
  \begin{Phonetics}{跑车}{pao3che1}[][HSK 7-9]
    \definition{s.}{bicicleta de corrida | \emph{roadster}; carro de corrida | carrinho para transportar toras em uma floresta}
    \definition{v.}{Coloquial: (condutores de trem) estar em serviço | (vagões de carvão em uma mina) deslizar acidentalmente para baixo (desgovernado) | Dialeto: dirigir um veículo de transporte | trabalhar em um trem; atendente de trem trabalhando no trem | escorregar acidentalmente para baixo; isso se refere a um acidente em um poço inclinado de mina, onde o cabo de aço se rompe repentinamente durante o içamento ou o guincho escorrega por outros motivos}
  \end{Phonetics}
\end{Entry}

\begin{Entry}{跑电}{12,5}{⾜,⽥}
  \begin{Phonetics}{跑电}{pao3dian4}
    \definition{s.}{fuga elétrica; quando o isolamento está danificado, a corrente escapa do fio ou aparelho para o exterior, o que também é chamado de corrente de fuga (漏电)}
  \synonymref{漏电}{lou4dian4}
  \end{Phonetics}
\end{Entry}

\begin{Entry}{跑龙套}{12,5,10}{⾜,⿓,⼤}
  \begin{Phonetics}{跑龙套}{pao3long2tao4}[][HSK 7-9]
    \definition{v.}{Teatro: interpretar um papel secundário | desempenhar um papel secundário; não ser ninguém | desempenhar um papel pequeno}
  \end{Phonetics}
\end{Entry}

\begin{Entry}{跑步}{12,7}{⾜,⽌}
  \begin{Phonetics}{跑步}{pao3//bu4}[][HSK 3]
    \definition{v.+compl.}{correr; trotar}
  \synonymref{奔跑}{ben1pao3}
  \antonymref{漫步}{man4bu4}
  \end{Phonetics}
\end{Entry}

\begin{Entry}{跑肚}{12,7}{⾜,⾁}
  \begin{Phonetics}{跑肚}{pao3 du4}
    \definition{v.}{Coloquial: ter diarréia}
  \synonymref{腹泻}{fu4xie4}
  \end{Phonetics}
\end{Entry}

\begin{Entry}{跑调}{12,10}{⾜,⾔}
  \begin{Phonetics}{跑调}{pao3diao4}
    \definition{v.}{Gíria: estar desafinado ou fora do tom (ao cantar)}
  \antonymref{好听}{hao3ting1}
  \antonymref{悦耳}{yue4'er3}
  \end{Phonetics}
\end{Entry}

\begin{Entry}{跑掉}{12,11}{⾜,⼿}
  \begin{Phonetics}{跑掉}{pao3diao4}
    \definition{v.}{fugir | dar meia"-volta}
  \antonymref{抓住}{zhua1 zhu4}
  \end{Phonetics}
\end{Entry}

\begin{Entry}{跑道}{12,12}{⾜,⾡}
  \begin{Phonetics}{跑道}{pao3dao4}[][HSK 7-9]
    \definition[条]{s.}{pista de decolagem; a pista de taxiagem utilizada pelas aeronaves durante a decolagem e o pouso | pista; pista de atletismo; as linhas brancas desenhadas na pista são usadas para corridas de corrida ou ciclismo}
  \end{Phonetics}
\end{Entry}

\begin{Entry}{跑腿}{12,13}{⾜,⾁}
  \begin{Phonetics}{跑腿}{pao3tui3}
    \definition{s.}{faz"-tudo, é uma figura que faz afazeres e pequenos serviços para as pessoas}
    \definition{v.}{fazer afazeres}
  \synonymref{代理}{dai4li3}
  \end{Phonetics}
\end{Entry}

\begin{Entry}{跑酷}{12,14}{⾜,⾣}
  \begin{Phonetics}{跑酷}{pao3ku4}
    \definition*{s.}{Empréstimo linguístico: Parkour}
  \end{Phonetics}
\end{Entry}

\begin{Entry}{跑题}{12,15}{⾜,⾴}
  \begin{Phonetics}{跑题}{pao3ti2}
    \definition{v.}{desviar"-se do assunto; afastar"-se do ponto principal; fazer uma tangente; tergiversar}
  \end{Phonetics}
\end{Entry}

%%%%%%%%%% 跟 %%%%%%%%%%
\subsection*{跟}\addcontentsline{loh}{figure}{跟}

\begin{Entry}{跟}{13}{⾜}
  \begin{Phonetics}{跟}{gen1}[][HSK 1]
    \definition{conj.}{e; expressa uma relação de união; 和}
    \definition{prep.}{com; Introduzir objetos relacionados à mesma ação, equivalente a 同 | para; em direção a | de; introduzir o objeto de comparação; equivalente a 从, 由 | como; objetos que causam comparações e semelhanças}
    \definition[个]{s.}{calcanhar; parte posterior do pé ou parte posterior do sapato ou meia | base (de um objeto)}
    \definition{v.}{seguir; acompanhar; seguir imediatamente na mesma direção | (uma mulher) estar casada com; casar"-se com alguém}
  \seealsoref{从}{cong2}
  \seealsoref{对}{dui4}
  \seealsoref{给}{gei3}
  \seealsoref{和}{he2}
  \seealsoref{随}{sui2}
  \seealsoref{同}{tong2}
  \seealsoref{由}{you2}
  \synonymref{嫁}{jia4}
  \synonymref{与}{yu3}
  \end{Phonetics}
\end{Entry}

\begin{Entry}{跟上}{13,3}{⾜,⼀}
  \begin{Phonetics}{跟上}{gen1shang5}[][HSK 7-9]
    \definition{v.}{acompanhar; alcançar; manter"-se a par de}
  \end{Phonetics}
\end{Entry}

\begin{Entry}{跟不上}{13,4,3}{⾜,⼀,⼀}
  \begin{Phonetics}{跟不上}{gen1bu5shang4}[][HSK 7-9]
    \definition{v.}{não é capaz de acompanhar; não conseguir alcançar}
  \end{Phonetics}
\end{Entry}

\begin{Entry}{跟前}{13,9}{⾜,⼑}
  \begin{Phonetics}{跟前}{gen1qian2}[][HSK 5]
    \definition{s.}{próximo; perto de; na frente de; (na ou para) a presença de alguém | o tempo imediatamente anterior a algum evento; tempo que se aproxima}
  \synonymref{附近}{fu4jin4}
  \synonymref{临近}{lin2jin4}
  \synonymref{旁边}{pang2bian1}
  \synonymref{身边}{shen1bian1}
  \synonymref{眼前}{yan3qian2}
  \synonymref{周围}{zhou1wei2}
  \synonymref{左右}{zuo3you4}
  \antonymref{后面}{hou4mian4}
  \antonymref{远处}{yuan3chu4}
  \antonymref{远方}{yuan3fang1}
  \end{Phonetics}
  \begin{Phonetics}{跟前}{gen1qian5}
    \definition{v.}{(filhos de alguém) viver com alguém (exclusivamente com relação à presença ou ausência de crianças)}
  \synonymref{旁边}{pang2bian1}
  \synonymref{身边}{shen1bian1}
  \synonymref{眼前}{yan3qian2}
  \synonymref{周边}{zhou1bian1}
  \synonymref{周围}{zhou1wei2}
  \synonymref{左右}{zuo3you4}
  \antonymref{后面}{hou4mian4}
  \end{Phonetics}
\end{Entry}

\begin{Entry}{跟随}{13,11}{⾜,⾩}
  \begin{Phonetics}{跟随}{gen1sui2}[][HSK 5]
    \definition{s.}{seguidor; usado para se referir a alguém que seguiu}
    \definition{v.}{seguir; ir atrás; acompanhar}
  \seealsoref{跟踪}{gen1zong1}
  \synonymref{伴随}{ban4sui2}
  \synonymref{陪同}{pei2tong2}
  \synonymref{追随}{zhui1sui2}
  \antonymref{带领}{dai4ling3}
  \antonymref{离开}{li2kai1}
  \antonymref{率领}{shuai4ling3}
  \antonymref{向导}{xiang4dao3}
  \end{Phonetics}
\end{Entry}

\begin{Entry}{跟踪}{13,15}{⾜,⾜}
  \begin{Phonetics}{跟踪}{gen1zong1}[][HSK 7-9]
    \definition{v.}{rastrear; alcançar; seguir; seguir atrás; seguir alguém; seguir os rastros de; seguir de perto}
  \end{Phonetics}
\end{Entry}

%%%%%%%%%% 跨 %%%%%%%%%%
\subsection*{跨}\addcontentsline{loh}{figure}{跨}

\begin{Entry}{跨}{13}{⾜}
  \begin{Phonetics}{跨}{kua4}[][HSK 6]
    \definition{adj.}{localizado ao lado de; anexo a}
    \definition{v.}{dar um passo; andar a passos largos | disputar; ficar de pernas abertas | atravessar; ir além (dos limites de uma certa quantidade, tempo, região, etc.)}
  \synonymref{骑}{qi2}
  \end{Phonetics}
\end{Entry}

\begin{Entry}{跨国}{13,8}{⾜,⼞}
  \begin{Phonetics}{跨国}{kua4guo2}[][HSK 7-9]
    \definition{adj.}{transnacional; que transcende fronteiras nacionais, envolvendo dois ou mais países}
  \end{Phonetics}
\end{Entry}

\begin{Entry}{跨越}{13,12}{⾜,⾛}
  \begin{Phonetics}{跨越}{kua4yue4}[][HSK 7-9]
    \definition{v.}{passar por cima; saltar por cima; atravessar; dar um passo largo; cruzar fronteiras regionais ou de período}
  \end{Phonetics}
\end{Entry}

%%%%%%%%%% 跪 %%%%%%%%%%
\subsection*{跪}\addcontentsline{loh}{figure}{跪}

\begin{Entry}{跪}{13}{⾜}
  \begin{Phonetics}{跪}{gui4}[][HSK 6]
    \definition{v.}{ajoelhar"-se; dobrar os joelhos de modo que um ou ambos os joelhos toquem o chão}
  \synonymref{拜}{bai4}
  \synonymref{坐}{zuo4}
  \antonymref{立}{li4}
  \end{Phonetics}
\end{Entry}

\begin{Entry}{跪拜}{13,9}{⾜,⼿}
  \begin{Phonetics}{跪拜}{gui4bai4}
    \definition{v.}{adorar de joelhos; prostrar"-se; curvar"-se}
  \synonymref{膜拜}{mo2bai4}
  \end{Phonetics}
\end{Entry}

%%%%%%%%%% 路 %%%%%%%%%%
\subsection*{路}\addcontentsline{loh}{figure}{路}

\begin{Entry}{路}{13}{⾜}
  \begin{Phonetics}{路}{lu4}[][HSK 1]
    \definition*{s.}{Sobrenome: Lu}
    \definition{clas.}{tipo; classe | linha; coluna; usado para um grupo de pessoas ou uma equipe; para organizar em ordem}
    \definition[条]{s.}{estrada; caminho; via | viagem; jornada; distância | maneira; meios | sequência; linha; lógica | região; distrito | rota | classe; classificação; grau | linha; fileira}
  \seealsoref{道路}{dao4lu4}
  \synonymref{道}{dao4}
  \synonymref{行}{hang2}
  \synonymref{类}{lei4}
  \synonymref{排}{pai2}
  \synonymref{途}{tu2}
  \end{Phonetics}
\end{Entry}

\begin{Entry}{路人}{13,2}{⾜,⼈}
  \begin{Phonetics}{路人}{lu4ren2}[][HSK 7-9]
    \definition{s.}{transeunte; estranho; o viajante se refere a uma pessoa sem qualquer relação com o caso}
  \end{Phonetics}
\end{Entry}

\begin{Entry}{路上}{13,3}{⾜,⼀}
  \begin{Phonetics}{路上}{lu4shang5}[][HSK 1]
    \definition{s.}{na estrada | a caminho; na rota; em processo de mudança de um lugar para outro}
  \synonymref{途中}{tu2zhong1}
  \end{Phonetics}
\end{Entry}

\begin{Entry}{路口}{13,3}{⾜,⼝}
  \begin{Phonetics}{路口}{lu4kou3}[][HSK 1]
    \definition[个]{s.}{cruzamento; intersecção; onde as estradas se encontram}
  \end{Phonetics}
\end{Entry}

\begin{Entry}{路子}{13,3}{⾜,⼦}
  \begin{Phonetics}{路子}{lu4zi5}[][HSK 7-9]
    \definition[个]{s.}{maneira; abordagem; meios | Coloquial: conexões sociais; atração}
  \end{Phonetics}
\end{Entry}

\begin{Entry}{路边}{13,5}{⾜,⾡}
  \begin{Phonetics}{路边}{lu4bian1}[][HSK 2]
    \definition{s.}{calçada; beira da estrada; margem da rua}
  \end{Phonetics}
\end{Entry}

\begin{Entry}{路灯}{13,6}{⾜,⽕}
  \begin{Phonetics}{路灯}{lu4deng1}[][HSK 7-9]
    \definition[盏]{s.}{poste de iluminação pública (ou de rua)}
  \end{Phonetics}
\end{Entry}

\begin{Entry}{路过}{13,6}{⾜,⾡}
  \begin{Phonetics}{路过}{lu4guo4}[][HSK 6]
    \definition{v.}{passar por (algum lugar); atravessar}
  \seealsoref{经过}{jing1guo4}
  \antonymref{停留}{ting2liu2}
  \end{Phonetics}
\end{Entry}

\begin{Entry}{路况}{13,7}{⾜,⼎}
  \begin{Phonetics}{路况}{lu4kuang4}[][HSK 7-9]
    \definition{s.}{condições da estrada (referentes à superfície da estrada, pavimento, fluxo de tráfego, etc.); tráfego}
  \end{Phonetics}
\end{Entry}

\begin{Entry}{路线}{13,8}{⾜,⽷}
  \begin{Phonetics}{路线}{lu4xian4}[][HSK 3]
    \definition[条]{s.}{rota; caminho; linha; a estrada percorrida de um lugar a outro | linha; diretriz (de política, ideologia, campo de trabalho); a via fundamental a seguir em termos ideológicos, políticos ou profissionais}
  \synonymref{道路}{dao4lu4}
  \synonymref{轨迹}{gui3ji4}
  \synonymref{门路}{men2lu5}
  \synonymref{途径}{tu2jing4}
  \synonymref{线路}{xian4lu4}
  \synonymref{准则}{zhun3ze2}
  \end{Phonetics}
\end{Entry}

\begin{Entry}{路段}{13,9}{⾜,⽎}
  \begin{Phonetics}{路段}{lu4duan4}[][HSK 7-9]
    \definition{s.}{trecho de uma rodovia ou ferrovia | trecho rodoviário; um trecho de estrada}
  \end{Phonetics}
\end{Entry}

\begin{Entry}{路面}{13,9}{⾜,⾯}
  \begin{Phonetics}{路面}{lu4mian4}[][HSK 7-9]
    \definition{s.}{superfície da estrada; pavimento}
  \end{Phonetics}
\end{Entry}

\begin{Entry}{路途}{13,10}{⾜,⾡}
  \begin{Phonetics}{路途}{lu4tu2}[][HSK 7-9]
    \definition{s.}{estrada; caminho | caminho; jornada}
  \end{Phonetics}
\end{Entry}

\begin{Entry}{路程}{13,12}{⾜,⽲}
  \begin{Phonetics}{路程}{lu4cheng2}[][HSK 7-9]
    \definition[段]{s.}{distância percorrida; viagem; descreve metaforicamente o processo de desenvolvimento de um evento | rota; caminho; distância; o comprimento total do percurso realizado por um objeto em movimento, desde seu ponto inicial até seu ponto final}
  \end{Phonetics}
\end{Entry}

%%%%%%%%%% 跳 %%%%%%%%%%
\subsection*{跳}\addcontentsline{loh}{figure}{跳}

\begin{Entry}{跳}{13}{⾜}
  \begin{Phonetics}{跳}{tiao4}[][HSK 3]
    \definition{v.}{pular; saltar | mover para cima e para baixo | pular (por cima); fazer omissões | quicar; a força elástica faz com que o objeto se mova repentinamente para cima | pulsar; palpitar; contrair"-se | pular sobre;  saltar sobre; cruzar}
  \seealsoref{蹦}{beng4}
  \synonymref{动}{dong4}
  \synonymref{闹}{nao4}
  \synonymref{弹}{tan2}
  \synonymref{腾}{teng2}
  \antonymref{静}{jing4}
  \end{Phonetics}
\end{Entry}

\begin{Entry}{跳水}{13,4}{⾜,⽔}
  \begin{Phonetics}{跳水}{tiao4shui3}[][HSK 6]
    \definition{s.}{Esporte: mergulho}
    \definition{v.}{mergulhar | Figurativo: (preços, lucros, etc.) cair drasticamente; cair repentinamente; mergulhar; despencar | cometer suicídio pulando na água | mergulhar (na água)}
  \antonymref{飙升}{biao1sheng1}
  \antonymref{上涨}{shang4zhang3}
  \end{Phonetics}
\end{Entry}

\begin{Entry}{跳电}{13,5}{⾜,⽥}
  \begin{Phonetics}{跳电}{tiao4dian4}
    \definition{v.}{(um disjuntor ou interruptor) desarmar}
  \end{Phonetics}
\end{Entry}

\begin{Entry}{跳伞}{13,6}{⾜,⼈}
  \begin{Phonetics}{跳伞}{tiao4//san3}[][HSK 7-9]
    \definition{s.}{salto de paraquedas}
    \definition{v.+compl.}{saltar de paraquedas; ejetar"-se}
  \end{Phonetics}
\end{Entry}

\begin{Entry}{跳动}{13,6}{⾜,⼒}
  \begin{Phonetics}{跳动}{tiao4dong4}[][HSK 7-9]
    \definition{v.}{bater; piscar; pular}
  \synonymref{跳跃}{tiao4yue4}
  \synonymref{震动}{zhen4dong4}
  \antonymref{静止}{jing4zhi3}
  \end{Phonetics}
\end{Entry}

\begin{Entry}{跳远}{13,7}{⾜,⾡}
  \begin{Phonetics}{跳远}{tiao4yuan3}[][HSK 3]
    \definition{s.}{salto em distância (atletismo)}
  \end{Phonetics}
\end{Entry}

\begin{Entry}{跳挡}{13,9}{⾜,⼿}
  \begin{Phonetics}{跳挡}{tiao4dang3}
    \definition{v.}{(um carro) escapar da marcha | sair da marcha}
  \end{Phonetics}
\end{Entry}

\begin{Entry}{跳蚤}{13,9}{⾜,⾍}
  \begin{Phonetics}{跳蚤}{tiao4zao5}
    \definition{s.}{pulga}
  \seealsoref{虼蚤}{ge4zao3}
  \end{Phonetics}
\end{Entry}

\begin{Entry}{跳高}{13,10}{⾜,⾼}
  \begin{Phonetics}{跳高}{tiao4gao1}[][HSK 3]
    \definition{s.}{salto em altura (atletismo)}
    \definition{v.}{saltar em altura}
  \seealsoref{跳高儿}{tiao4gao1r5}
  \synonymref{跳跃}{tiao4yue4}
  \end{Phonetics}
\end{Entry}

\begin{Entry}{跳高儿}{13,10,2}{⾜,⾼,⼉}
  \begin{Phonetics}{跳高儿}{tiao4gao1r5}
    \definition{s.}{Esporte: salto em altura}
  \end{Phonetics}
\end{Entry}

\begin{Entry}{跳绳}{13,11}{⾜,⽷}
  \begin{Phonetics}{跳绳}{tiao4sheng2}
    \definition{s.}{a corda usada para pular corda}
    \definition{v.}{pular corda; fazer pular corda; um esporte ou brincadeira infantil; girar uma corda em círculo; pular por cima da corda quando ela está perto do chão}
  \end{Phonetics}
\end{Entry}

\begin{Entry}{跳跃}{13,11}{⾜,⾜}
  \begin{Phonetics}{跳跃}{tiao4yue4}[][HSK 7-9]
    \definition{v.}{pular; saltar; dar um pulo}
  \synonymref{跨越}{kua4yue4}
  \synonymref{跳动}{tiao4dong4}
  \end{Phonetics}
\end{Entry}

\begin{Entry}{跳跳糖}{13,13,16}{⾜,⾜,⽶}
  \begin{Phonetics}{跳跳糖}{tiao4tiao4 tang2}
    \definition{s.}{\emph{Pop Rocks}, \emph{popping candy}}
  \end{Phonetics}
\end{Entry}

\begin{Entry}{跳频}{13,13}{⾜,⾴}
  \begin{Phonetics}{跳频}{tiao4pin2}
    \definition{s.}{FHSS, \emph{Frequency"-Hopping Spread Spectrum}, método de transmissão de sinais de rádio}
  \end{Phonetics}
\end{Entry}

\begin{Entry}{跳舞}{13,14}{⾜,⾇}
  \begin{Phonetics}{跳舞}{tiao4//wu3}[][HSK 3]
    \definition{v.+compl.}{dançar (como performance); executar dança, especialmente dança de salão}
  \synonymref{舞蹈}{wu3dao3}
  \antonymref{静止}{jing4zhi3}
  \antonymref{休息}{xiu1xi5}
  \end{Phonetics}
\end{Entry}

\begin{Entry}{跳槽}{13,15}{⾜,⽊}
  \begin{Phonetics}{跳槽}{tiao4//cao2}[][HSK 7-9]
    \definition{v.+compl.}{mudar de emprego; abandonar uma ocupação em favor de outra; trocar de emprego com frequência; mudar de emprego constantemente; essa metáfora descreve alguém que deixa seu emprego ou local de trabalho original para trabalhar em outra organização ou mudar de profissão | pular o cocho; o gado saiu do cocho que lhe fora designado para comer em outros cochos}
  \synonymref{离职}{li2//zhi2}
  \antonymref{坚守}{jian1shou3}
  \end{Phonetics}
\end{Entry}

%%%%%%%%%% 踊 %%%%%%%%%%
\subsection*{踊}\addcontentsline{loh}{figure}{踊}

\begin{Entry}{踊}{14}{⾜}
  \begin{Phonetics}{踊}{yong3}
    \definition{v.}{saltar; pular}[她兴奋地踊跃起来。===Ela pulou de alegria.]
  \end{Phonetics}
\end{Entry}

\begin{Entry}{踊跃}{14,11}{⾜,⾜}
  \begin{Phonetics}{踊跃}{yong3yue4}[][HSK 7-9]
    \definition{adj.}{positivo; entusiasmado; descreve um estado de grande entusiasmo e desejo de ser o primeiro}
    \definition{v.}{pular; saltar; dar um salto}
  \synonymref{奋勇}{fen4yong3}
  \synonymref{积极}{ji1ji2}
  \synonymref{主动}{zhu3dong4}
  \antonymref{消极}{xiao1ji2}
  \end{Phonetics}
\end{Entry}

%%%%%%%%%% 踏 %%%%%%%%%%
\subsection*{踏}\addcontentsline{loh}{figure}{踏}

\begin{Entry}{踏}{15}{⾜}
  \begin{Phonetics}{踏}{ta1}
    \definition{part.}{componente morfêmico}
  \end{Phonetics}
  \begin{Phonetics}{踏}{ta4}[][HSK 6]
    \definition{v.}{por os pés em; pisar em; esmagar com o pé | fazer uma investigação ou levantamento no local}
  \synonymref{踩}{cai3}
  \end{Phonetics}
\end{Entry}

\begin{Entry}{踏上}{15,3}{⾜,⼀}
  \begin{Phonetics}{踏上}{ta1shang5}[][HSK 7-9]
    \definition{v.}{pôr os pés em; pisar; pisar em | pisar em ou entrar; começar a entrar ou se envolver em uma área}
  \end{Phonetics}
\end{Entry}

\begin{Entry}{踏实}{15,8}{⾜,⼧}
  \begin{Phonetics}{踏实}{ta1shi5}[][HSK 6]
    \definition{adj.}{confiável; sério; estável e seguro; descreve uma atitude séria em relação ao trabalho ou estudo | à vontade; livre de ansiedade; descreve uma mente ou sentimento estável, sem qualquer preocupação ou ansiedade}
  \synonymref{安定}{an1ding4}
  \synonymref{安稳}{an1wen3}
  \synonymref{负责}{fu4ze2}
  \synonymref{实在}{shi2zai5}
  \synonymref{扎实}{zha1shi5}
  \antonymref{浮躁}{fu2zao4}
  \antonymref{冒险}{mao4//xian3}
  \end{Phonetics}
\end{Entry}

\begin{Entry}{踏板}{15,8}{⾜,⽊}
  \begin{Phonetics}{踏板}{ta4ban3}
    \definition{s.}{pedal; plataforma para os pés; apoio para os pés | banquinho (geralmente ao lado da cama) | pedal (em um carro, em um piano etc.) | trampolim (geralmente para salto em distância) | passarela; estrado para os pés | prancha de esteira}
  \end{Phonetics}
\end{Entry}

%%%%%%%%%% 踢 %%%%%%%%%%
\subsection*{踢}\addcontentsline{loh}{figure}{踢}

\begin{Entry}{踢}{15}{⾜}
  \begin{Phonetics}{踢}{ti1}[][HSK 6]
    \definition{v.}{chutar | jogar (por exemplo, futebol)}
  \synonymref{踹}{chuai4}
  \end{Phonetics}
\end{Entry}

\begin{Entry}{踢蹋舞}{15,17,14}{⾜,⾜,⾇}
  \begin{Phonetics}{踢蹋舞}{ti1 ta4 wu3}
    \definition{s.}{sapateado}
  \end{Phonetics}
\end{Entry}

\begin{Entry}{踢爆}{15,19}{⾜,⽕}
  \begin{Phonetics}{踢爆}{ti1 bao4}
    \definition{v.}{expor | revelar}
  \end{Phonetics}
\end{Entry}

%%%%%%%%%% 踩 %%%%%%%%%%
\subsection*{踩}\addcontentsline{loh}{figure}{踩}

\begin{Entry}{踩}{15}{⾜}
  \begin{Phonetics}{踩}{cai3}[][HSK 6]
    \definition{v.}{pisar; pisotear | pisar; metáfora: depreciar ou estragar | rastrear; antigamente significava rastrear (bandidos) ou investigar (casos)}
  \synonymref{贬}{bian3}
  \synonymref{踏}{ta4}
  \synonymref{抑}{yi4}
  \antonymref{吹}{chui1}
  \antonymref{捧}{peng3}
  \end{Phonetics}
\end{Entry}

%%%%%%%%%% 踹 %%%%%%%%%%
\subsection*{踹}\addcontentsline{loh}{figure}{踹}

\begin{Entry}{踹}{16}{⾜}
  \begin{Phonetics}{踹}{chuai4}[][HSK 7-9]
    \definition{v.}{chutar (com a sola do pé) | pisar; pisotear; pisar em}
  \end{Phonetics}
\end{Entry}

%%%%%%%%%% 蹦 %%%%%%%%%%
\subsection*{蹦}\addcontentsline{loh}{figure}{蹦}

\begin{Entry}{蹦}{18}{⾜}
  \begin{Phonetics}{蹦}{beng4}[][HSK 7-9]
    \definition{v.}{pular; saltar; quicar}
  \seealsoref{跳}{tiao4}
  \end{Phonetics}
\end{Entry}

\begin{Entry}{蹦极}{18,7}{⾜,⽊}
  \begin{Phonetics}{蹦极}{beng4ji2}
    \definition{s.}{\emph{bungee jumping}; um esporte radical em que a pessoa salta de uma altura enquanto está presa por uma corda elástica}
  \end{Phonetics}
\end{Entry}

%%%%%%%%%% 蹬 %%%%%%%%%%
\subsection*{蹬}\addcontentsline{loh}{figure}{蹬}

\begin{Entry}{蹬}{19}{⾜}
  \begin{Phonetics}{蹬}{deng1}[][HSK 7-9]
    \definition{v.}{pressionar com o pé; pisar; pisar em | Dialeto: calçar (sapatos ou calças); usar (sapatos) | Gíria: despejar (algo)}
  \end{Phonetics}
  \begin{Phonetics}{蹬}{deng4}
    \definition{s.}{lutar; ter dificuldade}
  \seealsoref{蹭蹬}{ceng4deng4}
  \end{Phonetics}
\end{Entry}

%%%%%%%%%% 蹭 %%%%%%%%%%
\subsection*{蹭}\addcontentsline{loh}{figure}{蹭}

\begin{Entry}{蹭}{19}{⾜}
  \begin{Phonetics}{蹭}{ceng4}[][HSK 7-9]
    \definition{v.}{esfregar; raspar; arranhar | esfregar em algo e ficar manchado; ser manchado com; manchar por fricção | mover"-se lentamente; demorar"-se; arrastar"-se | Dialeto: roubar}
  \end{Phonetics}
\end{Entry}

\begin{Entry}{蹭蹬}{19,19}{⾜,⾜}
  \begin{Phonetics}{蹭蹬}{ceng4deng4}
    \definition{interj.}{``Droga!''}
    \definition{v.}{enfrentar contratempos; estar sem sorte; ter má sorte}
  \end{Phonetics}
\end{Entry}

%%%%%%%%%% 蹲 %%%%%%%%%%
\subsection*{蹲}\addcontentsline{loh}{figure}{蹲}

\begin{Entry}{蹲}{19}{⾜}
  \begin{Phonetics}{蹲}{cun2}
    \definition{v.}{deslocar; torcer (uma perna); sofrer uma lesão articular}
  \end{Phonetics}
  \begin{Phonetics}{蹲}{dun1}[][HSK 6]
    \definition{v.}{agachamento sobre os calcanhares; dobrar as pernas o máximo possível, como se estivesse sentado, mas não deixar as nádegas tocarem o chão | ficar; metáfora para ficar ocioso em casa}
  \end{Phonetics}
\end{Entry}

\begin{Entry}{蹲下}{19,3}{⾜,⼀}
  \begin{Phonetics}{蹲下}{dun1xia4}
    \definition{v.}{agachar | agachar"-se}
  \end{Phonetics}
\end{Entry}

%%%%% EOF %%%%%

